\section{Model}
\label{sec:Model}

In this section, we introduce the model of the main experimental process that we study in this paper.
Our goal is to design pilot studies under a budget constraint that maximize impact---the expected improvement in outcomes---through an experimental process consisting of a pilot RCT, a potential follow-up RCT, and finally a treatment adoption decision, as described in the introduction and shown in~\Cref{fig:exp_process_flow}.

We consider a setting where, following each RCT stage, a statistical significance test serves as the \emph{only} continuation criterion for whether the intervention moves from one stage to the next.
As described in the previous section, this reflects a common continuation criterion used in practice.

We further assume that the pilot affects the outcome of the experimental process \emph{only} through its significance test, i.e.\ whether the treatment continues to a follow-up RCT.
In other words, the potential follow-up design is fixed, e.g.\ to something like a representative sample, and the adoption decision only depends on the outcomes of the follow-up RCT.
This assumption is necessary for tractability since designing the follow-up is usually itself a non-convex optimization problem.

The remainder of this section gives a formal definition of the model and objective.

\subsection{Three-stage experiment pipeline}

We model the experimental process in terms of three stages: the pilot RCT, the follow-up RCT, and adoption, indexed by $t\in\{1,2,3\}$, respectively.

In the pilot and follow-up RCT stages, a sample of subjects is recruited into the study.
At the adoption stage, the intervention is rolled out to a predetermined target population.
Each stage has a corresponding \emph{average treatment effect},
\[
  \ATE_t\in\R,\quad t\in\{1,2,3\},
\]
and the two RCT stages each have an \emph{average treatment effect estimate} derived from the RCT sample,
\[
  \widehat{\ATE}_t\in\R,\quad t\in\{1,2\}.
\]

Following each RCT stage $t=1,2$, the intervention moves to the next stage if and only if the average treatment effect estimate clears a threshold given by a \emph{statistical significance test},
\[
  \Test_t\in\{0,1\},\quad t\in\{1,2\},
\]
where we write $\Test_t=1$ if the treatment continues to the next stage and $\Test_t=0$ otherwise.
Since a treatment is only adopted if it passes both significance tests, we can write the \emph{realized improvement in average outcomes} as
\[
  \Test_1 \cdot \Test_2 \cdot\, \ATE_3.
\]

\subsection{RCT samples and the target population}

Each stage of the experimental process has a sample or population of units that can vary in its composition.
We define discrete unit types to capture feature variability, and the composition of each stage in terms of the number of units of each type.

\paragraph{Unit types}
There are $D\ge 2$ types of units (e.g., patients, users, schools, etc.) that can be treated or observed, indexed by $d\in\{1,\ldots,D\}$.
These types can vary in their treatment effects and sampling costs.
For example, in an education study where schools are the units, a type could include schools with similar demographics, funding, and size---all of which could affect the treatment effect---and at a similar distance from the research institution running the pilot study---which could affect the cost of sampling.

\paragraph{Pilot RCT sample}
At the pilot stage $t=1$, the pilot designer chooses how many units of each type to include in the study.
Let $s_{1d}\geq 0$ denote the number of units sampled from each type $d$, and $\bm{s}_1\in\R_{\geq 0}^D$ the vector of units sampled across all $D$ types.
Let $n_1 \defeq \bm{1}^{\top}\bm{s}_1$ be the pilot sample size.
For tractability, we let $\bm{s}_1$ take on continuous values, rather than imposing an integer constraint.
We further constrain the sample to be non-empty, i.e., $\bm{s}_1 \neq \bm{0}$ and $n_1 > 0$, since the average treatment effect is undefined for an empty sample, and write this as $\bm{s}_1 \in \R_{\geq 0}^{D} \setminus \{\bm{0}\}$.

\paragraph{Follow-up RCT sample}
At the follow-up stage $t=2$, a sample of units is enrolled in the follow-up RCT.
Let $\bm{s}_2\in\R_{\geq 0}^{D} \setminus \{\bm{0}\}$ be the vector of units sampled across all $D$ types if the follow-up is run, and $n_2 \defeq \bm{1}^{\top}\bm{s}_2 > 0$ the total follow-up sample size.
We focus on designing the optimal pilot sample $\bm{s}_1$ for a fixed follow-up study sample $\bm{s}_2$, and leave their optimal joint design for future work.

\paragraph{Target population}
At the adoption stage $t=3$, a fixed set of units we call the target population is treated if the treatment is adopted.
Let $\bm{s}_3\in\R_{\geq 0}^{D} \setminus \{\bm{0}\}$ be the vector of population sizes for each type, and $n_3 \defeq \bm{1}^{\top}\bm{s}_3 > 0$ be the total target population size.

\begin{figure}[ht]
  \centering
  
\newcommand{\drawdownarc}[4][]{%
  \path let \p1=(#2), \p2=(#3) in
  \pgfextra{
    \pgfmathsetmacro{\dx}{\x2-\x1}
    \pgfmathsetmacro{\dy}{\y2-\y1}
    \pgfmathsetmacro{\L}{sqrt(\dx*\dx+\dy*\dy)}

    \pgfmathsetlengthmacro{\slen}{#4cm}
    \pgfmathsetmacro{\s}{\slen}

    \pgfmathsetmacro{\r}{(\L*\L)/(8*\s) + (\s/2)}
    \pgfmathsetmacro{\d}{\r-\s}

    \pgfmathsetmacro{\xm}{(\x1+\x2)/2}
    \pgfmathsetmacro{\ym}{(\y1+\y2)/2}

    \pgfmathsetmacro{\nx}{\dy/\L}
    \pgfmathsetmacro{\ny}{-\dx/\L}
    \pgfmathsetmacro{\flip}{ifthenelse(\ny<0,-1,1)}
    \pgfmathsetmacro{\nx}{\flip*\nx}
    \pgfmathsetmacro{\ny}{\flip*\ny}

    \pgfmathsetmacro{\xc}{\xm + \nx*\d}
    \pgfmathsetmacro{\yc}{\ym + \ny*\d}

    \pgfmathsetmacro{\angA}{atan2(\y1-\yc,\x1-\xc)}
    \pgfmathsetmacro{\theta}{2*asin(\L/(2*\r))}

    \pgfmathsetmacro{\sgn}{ifthenelse(\dx>0,1,-1)}
    \pgfmathsetmacro{\delta}{\sgn*\theta}

    \xdef\ArcStart{\angA}
    \xdef\ArcDelta{\delta}
    \xdef\ArcRadius{\r}
  };
  \draw[#1] (#2) arc[start angle=\ArcStart, delta angle=\ArcDelta,
  radius=\ArcRadius pt];
}

\begin{tikzpicture}[
    font=\small,
    >=Triangle,
    stage/.style={align=center},
    label/.style={align=center, font=\bfseries},
    flow/.style={->, line width=0.75pt, line cap=round},
    decisionflow/.style={flow, dashed}
  ]

  \def\colsep{3.5}
  \def\toplabsep{0.40}
  \def\rowsep{1.65}
  \def\arrowgap{0.15}

  \def\curvedrop{0.40}
  \def\curvegap{0.18}
  \def\midgap{0.18}

  \def\testdrop{1.35}

  \def\boxpadx{9pt}
  \def\boxpady{7pt}

  \coordinate (c1) at (0,0);
  \coordinate (c2) at (\colsep,0);
  \coordinate (c3) at (2*\colsep,0);

  \node[stage] (S1) at (c1) {Pilot Sample $\bm{s}_1$};
  \node[stage] (S2) at (c2) {Follow-up Sample $\bm{s}_2$};
  \node[stage] (S3) at (c3) {Population $\bm{s}_3$};

  \node[label] (DecLbl)     at ($(S1.north)+(0,\toplabsep)$) {Decision};
  \node[label] (FixTopLbl)  at
  ($($(S2.north)!0.5!(S3.north)$)+(0,\toplabsep)$) {Fixed};

  \node[stage] (H1) at ($(S1)+(0,-\rowsep)$) {$\widehat{\ATE}_1$};
  \node[stage] (H2) at ($(S2)+(0,-\rowsep)$) {$\widehat{\ATE}_2$};
  \node[stage] (Y)  at ($(S3)+(0,-\rowsep)$) {$\ATE_3$};

  \node[stage, text=red] (T1) at ($($(H1)!0.5!(H2)$)+(0,-\testdrop)$)
  {$\Test_1$};
  \node[stage, text=red] (T2) at ($($(H2)!0.5!(Y)$ )+(0,-\testdrop)$)
  {$\Test_2$};

  \begin{pgfonlayer}{background}
    \node (BoxDec) [
      fit=(DecLbl)(S1),
      fill=blue!25,
      inner xsep=\boxpadx,
      inner ysep=\boxpady
    ] {};

    \node (BoxFixTop) [
      fit=(FixTopLbl)(S2)(S3),
      fill=gray!25,
      inner xsep=\boxpadx,
      inner ysep=\boxpady
    ] {};

  \end{pgfonlayer}

  \coordinate (S2BoxSouth) at (S2.south |- BoxFixTop.south);
  \coordinate (S3BoxSouth) at (S3.south |- BoxFixTop.south);

  \draw[flow] ($(BoxDec.south)+(0,-\arrowgap)$)   --
  ($(H1.north)+(0,\arrowgap)$);
  \draw[flow] ($(S2BoxSouth)+(0,-\arrowgap)$)    --
  ($(H2.north)+(0,\arrowgap)$);
  \draw[flow] ($(S3BoxSouth)+(0,-\arrowgap)$)    -- ($(Y.north)+(0,\arrowgap)$);

  \coordinate (ArcAstart) at ($(H1.south)+(0,-\curvegap)$);
  \coordinate (ArcAend)   at ($(H2.south)+(-\midgap,-\curvegap)$);
  \coordinate (ArcBstart) at ($(H2.south)+(\midgap,-\curvegap)$);
  \coordinate (ArcBend)   at ($(Y.south)+(0,-\curvegap)$);

  \drawdownarc[decisionflow]{ArcAstart}{ArcAend}{\curvedrop}
  \drawdownarc[decisionflow]{ArcBstart}{ArcBend}{\curvedrop}

\end{tikzpicture}

  \caption{
    \textbf{The three-stage experimental pipeline.}
    The pilot design $\bm{s}_1$ is highlighted in blue.
    The samples $\bm{s}_1$ and $\bm{s}_2$ in the first two stages determine the distribution of the average treatment effect estimates, and the population $\bm{s}_3$ in the third stage determines the change in outcomes if both significance tests pass.
  }
  \label{fig:experiment_process_flow}
\end{figure}

\subsection{Probabilistic model}

We define a probabilistic model of the experimental process.
The pilot designer has a prior over heterogeneous average treatment effects for each type.
Under this prior, the sample compositions of the pilot and follow-up RCTs and the composition of the target population induce a distribution over their average treatment effects.
Combined with a sampling distribution for sample average treatment effect estimates and a definition of the significance tests, this gives a full probability model for the experimental process.

We now define the heterogeneous treatment effect vector and how the sample and population compositions determine the distribution of average and sample average treatment effects.

\paragraph{Heterogeneous treatment effects}
We denote the conditional average treatment effect for type $d$ as $\tau_d$ and
let $\bm{\tau}\in\R^D$ be the vector of type average treatment effects.%
\footnote{
  Under the potential outcomes framework, these can be written as the mean
  difference in outcomes under treatment $Y_{i}(1)$ and control $Y_{i}(0)$ for a
  unit $i$ of type $d$, $\tau_d = \E[Y_{i}(1) - Y_{i}(0) \mid i \text{ is of
  type } d]$.
}
The experimenter has a prior over the conditional average treatment effects across types.
Our main results hold under general priors $\mathcal{P}$ on $\R^D$ satisfying moment and other mild regularity conditions,
\[
  \bm{\tau}\sim\mathcal{P}.
\]
The existence of fourth moments of $\bm\tau$ and a Lebesgue density positive on an open ball containing the origin are sufficient for all the main results; the appendix gives specific conditions for each.

We are able to give a more interpretable characterization of the optimal small-budget pilot design when further restricting the prior to the class of elliptical distributions~\cite{FangKo90}
\[
  \bm{\tau}\sim\mathrm{EC}_D(\bm{\mu},\Sigma,\phi),
\]
which includes and can be viewed as generalizing the multivariate normal distribution.

\paragraph{Average Treatment Effects}
Each stage $t=1,2,3$ and its associated sample $\bm{s}_t$ has an average treatment effect $\ATE_t$ given in terms of the type-level effects as
\[
  \ATE_t \defeq \frac{\bm{s}_t^{\top}\bm{\tau}}{n_t},
  \quad t\in\{1,2,3\}.
\]
These effects are not known to the pilot designer, but their prior distributions are induced by the designer's prior distribution on $\bm\tau$.

\paragraph{RCT Average Treatment Effect Estimates}
In the pilot and follow-up, the experimenter employs an RCT design paired with an appropriate estimator of the average treatment effect.
We approximate the sampling distributions of the estimators using independent normal distributions across stages, with precision proportional to the sample size $n_t$.%
\footnote{
  We set the variance to be exactly equal to the inverse sample size for notational simplicity; one could achieve this by scaling the outcome by a design- and estimator-dependent constant.
  If each type has different sampling noise, one can interpret the costs as costs-per-precision instead of costs-per-unit.
}
This approximation is motivated by central limit theorems that apply to many standard designs and estimators, such as completely- and Bernoulli-randomized designs as well as matched pairs and more general stratified designs~\cite{LiDi17,Ding23}, and is frequently used when modeling sampling noise~\cite{FrazierPo09,ChickGa22,AzevedoDe20}.
\[
  \widehat{\ATE}_t \mid \bm{\tau} \overset{\text{ind}}{\sim}
  \mathcal{N}\left(
    \ATE_t,
    \frac{1}{n_t}
  \right),
  \quad t\in\{1,2\}.
\]

\paragraph{Significance Tests as Continuation Criteria}
Following each RCT stage, the treatment advances to the next stage if the sample average treatment effect clears a threshold given by a statistical test with a specified significance level.

Formally, the continuation criteria are two Z-tests with one-sided%
\footnote{
  The choice of a one-sided test is without loss of generality, as in this case
  a one-sided test at level $\alpha$ is equivalent to a two-sided test at level
  $2\alpha$ when continuation requires a positive significant result.
}
significance levels $\alpha_{1}\in(0,1)$ and $\alpha_{2}\in(0,1)$.
Under the weak null that the treatment has no true average effect $\ATE_t=0$, the following test statistic has a standard normal sampling distribution by the distribution of the average treatment effect estimate:
\[
  \sqrt{n_t}\cdot\widehat{\ATE}_t
  =
  \sqrt{n_t}\cdot(\widehat{\ATE}_t - \ATE_t)
  \sim
  \mathcal{N}(0,1),
  \quad t\in\{1,2\}.
\]
Thus, one-sided tests at the $\alpha_{1}$ and $\alpha_{2}$ significance levels that allow continuation only for significant positive effects are given by the following decision rules,
\begin{align*}
  \Test_t &
  \defeq
  \indic{
    \sqrt{n_t}\cdot\widehat{\ATE}_t\ge
    z_{1-\alpha_t}
  } ,\quad z_{1-\alpha_t}\defeq\Phi^{-1}\left(1-\alpha_{t}\right)
\end{align*}
where we write $\Phi^{-1}$ for the normal inverse cumulative distribution function and $z_{1-\alpha_t}$ for the corresponding critical value of the test statistic.

\subsection{Pilot constraint, impact objective, and optimization problem}
\label{sec:pilot-budget-constraint-and-objective}

\paragraph{Pilot Sampling Costs}
For each type $d$, it costs $c_{d}>0$ to recruit a unit of type $d$ into the pilot RCT.
This cost can represent a monetary cost in enrolling the unit, a time cost in recruiting the unit or measuring its outcomes, etc.

\paragraph{Pilot Budget and Investment}
The pilot designer has a fixed budget $b > 0$ available to invest in running the pilot RCT.
When choosing the pilot design $\bm{s}_1$, the experimenter is subject to the constraint that the investment $I \defeq \bm{c}^{\top}\bm{s}_{1}$ is bounded by the budget $b$,
\[
  \bm{c}^{\top}\bm{s}_{1}\le b.
\]
We assume that the pilot budget constraint is the only binding constraint on the pilot sample.

\paragraph{Pilot RCT Impact}
Our main goal is to understand how to optimize the expected realized improvement in
outcomes in the target population---what we call the \emph{pilot RCT impact}.
The true average treatment effect is only realized in the target population if
the trial continues to a follow-up and results in adoption, so we can define the
pilot impact as%
\footnote{
  Note that writing $\ATE_3$ in the expression for the pilot impact, i.e.\ ignoring any finite sample variation in the target population average treatment effect, is without loss of generality since we are interested in the expected value.
  We could substitute some $\widehat{\ATE}_3$ for $\ATE_3$ in the above expression, where so long as $\E[\widehat{\ATE}_3\mid \bm{\tau}] = \ATE_3$ and the noise is independent given $\bm{\tau}$, the expected pilot impact would be unchanged.
}
\[
  V(\bm{s}_1) \defeq \E\left[
    \Test_1
    \cdot
    \Test_2
    \cdot\,
    \ATE_3
  \right]
\]
Other objectives can be formulated within our framework, such as maximizing the probability of passing the pilot statistical test $\E[\Test_1]$ or a sequence of tests $\E[\Test_1 \cdot \Test_2]$, or the impact of making an adoption decision based on a single RCT $\E[\Test_1 \cdot \ATE_3]$.
We compare small-budget optimal policies for these alternate objectives in \Cref{sec:alternate-objectives-appendix}.

We can formalize our mathematical goal as the following optimization problem.

\vspace{0.1cm}
\medskip
\noindent\textbf{Pilot RCT Impact Maximization Problem}
\vspace{0.1cm}
\begin{align*}
  \text{maximize}\quad  &
  V(\bm{s}_1)\\
  \text{s.t.}\quad &
  \begin{aligned}
    \bm{c}^{\top}\bm{s}_{1} & \le b &  & \text{(pilot
    budget constraint)}\\
    \bm{s}_{1} & \in\R_{\geq 0}^{D} \setminus \{\bm{0}\} &  &
    \text{(non-negative, non-zero
    sample)}
  \end{aligned}
\end{align*}

\vspace{0.2cm}
\noindent We note that in general this problem is non-convex, even under the assumption of a Gaussian prior on treatment effects that is commonly made in the literature on clinical trials and R\&D experiments in operations research~\cite{FrazierPo09,ChickGa22}.
For this reason, we will focus on the set of optimal solutions
\[
  \mathcal{S}_1^{\star}
  =
  \underset{\bm{s}_1}{\arg\max}
  \left\{
    V(\bm{s}_1) :
    \bm{c}^\top\bm{s}_1 \leq b,\
    \bm{s}_1 \in \R_{\geq 0}^{D} \setminus \{\bm{0}\}
  \right\}.
\]

In the following sections, we characterize this set of optimal solutions in two parameter regimes: for  large pilot budgets ($b\to\infty$) in \Cref{sec:large-budget} and small pilot budgets ($b$ smaller than a finite threshold) in \Cref{sec:small-budget}.
We write $\mathcal{S}_1^{\star}(b)$ when we wish to emphasize the dependence of the set of optimal solutions on the pilot budget $b$.
To start, we note that an optimal solution always exists.

\begin{proposition}[Existence of optimal
  solution]
  \label{prop:optimal-solution-exists}
  Suppose $\E[|\ATE_3|]<\infty$ and $\var(\ATE_3)>0$.
  Then the impact maximization problem has at least one optimal solution, i.e.\ $\mathcal{S}_1^{\star}\neq\emptyset$.
\end{proposition}

This follows from the observation that the objective is continuous and does not approach its supremum as $\bm{s}_1\to\bm{0}$, where $\bm{0}$ is the only limit point not in the feasible set since every valid design must have a positive sample size so that the objective is well-defined.
See \Cref{\sectionof{prop:optimal-solution-exists-appendix}} for the
\hyperref[prop:optimal-solution-exists-appendix]{proof}.

We will frequently reparametrize the decision variable $\bm{s}_1$ in terms of the pilot investment $\invst=\bm{c}^{\top}\bm{s}_1\in\R$ and a unit-investment normalized sample
\[
  \overline{\bm{s}}_1\in\{\bm{x}\in\R_{\geq 0}^D : \bm{c}^\top
  \bm{x}=1\}\eqdef\overline{\mathcal{S}}.
\]
By the definition of the feasible set above,
\[
  \bm{s}_1 = \invst\,\overline{\bm{s}}_1\ \ \text{is feasible if and
  only if}\ \ 0 < \invst \le
  b\ \ \text{and}\ \ \overline{\bm{s}}_1\in\overline{\mathcal{S}}.
\]
